%! Matrix Multiplication and A = CR — Unit 1, Lesson 1.4 — Guided Notes (Answer Key)
\documentclass[10pt]{article}
\usepackage{linalg-article}
\usepackage{linalg-key}   % pulls in -boxes; do NOT also load -boxes
\newcommand{\TallMath}[1]{$\displaystyle #1\rule[-1.4em]{0pt}{3.2em}$}
\newcommand{\xx}{\mathbf{x}}
\newcommand{\bb}{\mathbf{b}}

\begin{document}

\pageheader{Unit 1, Lesson 1.4}{Guided Notes --- Answer Key}
\namedateperiod

\vspace{0.12in}

\begin{objectivebox}
\textbf{By the end of this lesson, I will be able to\dots}
\begin{itemize}\setlength{\itemsep}{2pt}
  \item multiply matrices \emph{by columns}: each column of $AB$ is $A$ times a column of $B$;
  \item factor a matrix as $A=CR$ --- pull the \emph{independent} columns into $C$ and build the \emph{recipe} $R$;
  \item state the \emph{rank} (the number of independent columns) and explain what it means.
\end{itemize}
\end{objectivebox}

\vspace{0.10in}

\begin{vocabbox}
\small
Fill in each term as we build it together.
\vspace{2pt}
\vocabans{Matrix multiplication (by columns)}{column $j$ of $AB$ is $A\mathbf{b}_j$; a combination of $A$'s columns}
\vocabans{Independent columns / $C$}{the columns that add a new direction; collected in $C$ (a basis for $C(A)$)}
\vocabans{Recipe $R$ (rebuilds each column)}{each column of $R$ gives the weights that rebuild a column of $A$ from $C$; $A=CR$}
\vocabans{Rank}{the number of independent columns --- the number of columns in $C$}
\end{vocabbox}

\begin{hookbox}
Look at the three columns of this matrix --- is any one just a mix of the others?
\[
  A=\begin{bmatrix} 1 & 3 & 4 \\ 2 & 1 & 3 \end{bmatrix},
  \qquad
  \mathbf{a}_1=\begin{bmatrix} 1 \\ 2 \end{bmatrix},\ \
  \mathbf{a}_2=\begin{bmatrix} 3 \\ 1 \end{bmatrix},\ \
  \mathbf{a}_3=\begin{bmatrix} 4 \\ 3 \end{bmatrix}.
\]
Is $\mathbf{a}_3$ a combination of $\mathbf{a}_1$ and $\mathbf{a}_2$? Try
$\mathbf{a}_1+\mathbf{a}_2=\ans{$\begin{bsmallmatrix} 4\\3 \end{bsmallmatrix}$}$. So only \ans{two}
columns are really new. \ansline{$\mathbf{a}_3=\mathbf{a}_1+\mathbf{a}_2$, so it is a repeat --- only $\mathbf{a}_1,\mathbf{a}_2$ are independent.}
\end{hookbox}

\begin{notesbox}{1. Multiplying Matrices --- One Column at a Time}
To multiply $AB$, use \emph{each column of $B$} to combine the columns of $A$. Column $j$ of the
product is $A$ times column $j$ of $B$:
\[
  AB = A\begin{bmatrix} \mathbf{b}_1 & \mathbf{b}_2 \end{bmatrix}
  = \begin{bmatrix} A\mathbf{b}_1 & A\mathbf{b}_2 \end{bmatrix}.
\]
\textit{Example.} $A=\begin{bmatrix} 1 & 3 \\ 2 & 1 \end{bmatrix}$,
$B=\begin{bmatrix} 2 & 1 \\ 1 & 0 \end{bmatrix}$:
\[
  A\mathbf{b}_1 = 2\begin{bmatrix} 1\\2 \end{bmatrix}+1\begin{bmatrix} 3\\1 \end{bmatrix}=\begin{bmatrix} 5\\5 \end{bmatrix},
  \quad
  A\mathbf{b}_2 = 1\begin{bmatrix} 1\\2 \end{bmatrix}+0\begin{bmatrix} 3\\1 \end{bmatrix}={\color{keyred}\begin{bsmallmatrix} \mathbf{1}\\\mathbf{2} \end{bsmallmatrix}},
  \quad AB={\color{keyred}\begin{bsmallmatrix} \mathbf{5}&\mathbf{1}\\\mathbf{5}&\mathbf{2} \end{bsmallmatrix}}.
\]
Every column of a product is a \textbf{combination of $A$'s columns} --- Lesson~1.3, once per column.
\end{notesbox}

\begin{notesbox}{2. The Independent Columns: $C$}
Some columns add a \emph{new direction}; others are just repeats (combinations) of ones you already
have. Collect the \textbf{independent} columns --- the truly new ones --- into a matrix $C$.
\begin{center}
\begin{tikzpicture}[scale=0.62]
  \draw[step=1, linegray!60, thin] (-0.3,-0.3) grid (4.3,3.3);
  \draw[->, thick, charcoal] (-0.4,0)--(4.5,0) node[right]{\scriptsize $x$};
  \draw[->, thick, charcoal] (0,-0.4)--(0,3.5) node[above]{\scriptsize $y$};
  \draw[->, very thick, burgundy] (0,0)--(1,2) node[above left=-2pt]{\scriptsize $\mathbf{a}_1$};
  \draw[->, very thick, royalblue] (0,0)--(3,1) node[below right=-2pt]{\scriptsize $\mathbf{a}_2$};
  \draw[->, very thick, black!70, dashed] (0,0)--(4,3) node[above right=-2pt]{\scriptsize $\mathbf{a}_3=\mathbf{a}_1+\mathbf{a}_2$};
  \node[charcoal] at (2.0,-1.05) {\scriptsize $\mathbf{a}_3$ is a repeat: only $\mathbf{a}_1,\mathbf{a}_2$ are independent};
\end{tikzpicture}
\end{center}
For $A=\begin{bmatrix} 1 & 3 & 4 \\ 2 & 1 & 3 \end{bmatrix}$: columns~1 and~2 are not parallel, and
column~3 is their sum. The independent columns are \ans{$\begin{bsmallmatrix} 1\\2 \end{bsmallmatrix}$}
and \ans{$\begin{bsmallmatrix} 3\\1 \end{bsmallmatrix}$}, so
\[
  C={\color{keyred}\begin{bmatrix} \mathbf{1} & \mathbf{3} \\[3pt] \mathbf{2} & \mathbf{1} \end{bmatrix}}.
\]
\end{notesbox}

\begin{notesbox}{3. The Recipe $R$, and $A=CR$}
$R$ is the \textbf{recipe card}: each column of $R$ says how to mix the columns of $C$ to rebuild
that column of $A$. For a column that \emph{is} in $C$, the recipe is just ``take that one'' (a
column of the identity). For a repeat, the recipe is its weights.
\[
  \text{col 1}=1\,\mathbf{c}_1+0\,\mathbf{c}_2\to\begin{bmatrix} 1\\0 \end{bmatrix},\quad
  \text{col 2}\to\begin{bmatrix} 0\\1 \end{bmatrix},\quad
  \text{col 3}=\mathbf{c}_1+\mathbf{c}_2\to\begin{bmatrix} 1\\1 \end{bmatrix}.
\]
So $R=\begin{bmatrix} 1 & 0 & 1 \\ 0 & 1 & 1 \end{bmatrix}$ and $A=CR$. Check by multiplying by
columns: $C\begin{bsmallmatrix} 1\\1 \end{bsmallmatrix}=\mathbf{c}_1+\mathbf{c}_2=\ans{$\begin{bsmallmatrix} 4\\3 \end{bsmallmatrix}$}$
(matches $\mathbf{a}_3$? \ansline{Yes --- $\begin{bsmallmatrix} 4\\3 \end{bsmallmatrix}=\mathbf{a}_3$, so $CR$ rebuilds column~3.})
\end{notesbox}

\begin{notesbox}{4. Rank: How Many Columns Are Really New}
The \textbf{rank} is the number of independent columns --- the number of columns in $C$. It counts
how many directions the columns truly span.
\begin{itemize}\setlength{\itemsep}{1pt}
  \item Independent columns fill the whole \textbf{plane} $\Rightarrow$ rank \ans{2}.
  \item Parallel columns fill only a \textbf{line} $\Rightarrow$ rank \ans{1}.
\end{itemize}
For the $A$ above, $C$ has two columns, so $\operatorname{rank}(A)=\ans{2}$.
Your turn: $B=\begin{bmatrix} 1 & 2 \\ 2 & 4 \end{bmatrix}$ has parallel columns --- what is its
rank, and what is $C$? \ansline{Rank $1$; $C=\begin{bsmallmatrix} 1\\2 \end{bsmallmatrix}$ (one column), and $R=\begin{bsmallmatrix} 1 & 2 \end{bsmallmatrix}$.}
\end{notesbox}

\begin{practicebox}
\begin{enumerate}[leftmargin=1.6em]
  \item Multiply by columns: $\begin{bmatrix} 1 & 0 \\ 3 & 1 \end{bmatrix}\begin{bmatrix} 2 & 1 \\ 1 & 4 \end{bmatrix}$. \hfill \ans{$\begin{bsmallmatrix} 2\\7 \end{bsmallmatrix}$ and $\begin{bsmallmatrix} 1\\7 \end{bsmallmatrix}$, so $\begin{bsmallmatrix} 2 & 1\\7 & 7 \end{bsmallmatrix}$}
  \item Factor $A=\begin{bmatrix} 1 & 2 & 3 \\ 0 & 1 & 1 \end{bmatrix}$ as $A=CR$. Which columns are
        independent? Write $C$ and $R$, and give the rank. \ansline{Columns~1 and~2 are independent; column~3 $=$ col~1 $+$ col~2. $C=\begin{bsmallmatrix} 1 & 2\\0 & 1 \end{bsmallmatrix}$, $R=\begin{bsmallmatrix} 1 & 0 & 1\\0 & 1 & 1 \end{bsmallmatrix}$, rank $2$.}
  \item True or false, with a reason: a $2\times 4$ matrix can have rank $3$. \ansline{False --- each column is a vector in the plane ($2$ entries), so at most $2$ can be independent; rank $\le 2$.}
\end{enumerate}
\end{practicebox}

\end{document}
